% !TeX program = xelatex
% 独立可编译源文件；参考文献与技术路线图均在本文件内。
% 正文按申请书模板的一、二、三部分组织。题目另填模板“项目名称”。
% 申请金额按重点课题上限30万元测算；技术指标为拟定目标。
% 已有工作依据：26-5356_01_MS.pdf、26-0912.pdf及用户指定的三个代码目录。
% 两篇匿名稿件不列作已发表论文；机器人代码仅作只读核查。
\documentclass[UTF8,a4paper,zihao=-4,fontset=fandol]{ctexart}
\usepackage[left=2.6cm,right=2.6cm,top=2.5cm,bottom=2.5cm]{geometry}
\usepackage{amsmath,amssymb,bm,booktabs,longtable,array,tabularx}
\usepackage{enumitem,setspace,xcolor,caption,needspace,placeins}
\usepackage{tikz}
\usetikzlibrary{arrows.meta,positioning,calc}
\usepackage[unicode,hidelinks]{hyperref}
\usepackage{xurl}
\hypersetup{pdftitle={柔性支撑空间机器人耦合动力学与学习增强阻抗控制},pdfsubject={南航重点实验室开放课题申请书报告正文}}
\setstretch{1.32}
\setlength{\parindent}{2em}
\setlength{\parskip}{0.15em}
\setlength{\emergencystretch}{2em}
\clubpenalty=10000
\widowpenalty=10000
\displaywidowpenalty=10000
\raggedbottom
\setlist[enumerate]{leftmargin=2.2em,itemsep=0.3em,topsep=0.3em}
\captionsetup{font=small,labelfont=bf,labelsep=quad,skip=8pt}
\ctexset{
 section={format=\zihao{4}\bfseries,name={,、},number=\chinese{section},aftername={},beforeskip=1.2em,afterskip=0.7em},
 subsection={format=\zihao{-4}\bfseries,number=\arabic{subsection},aftername={.\ },beforeskip=1em,afterskip=0.5em},
 subsubsection={format=\zihao{-4}\bfseries,number=\arabic{subsection}.\arabic{subsubsection},aftername={\quad},beforeskip=0.8em,afterskip=0.3em}
}
\newcommand{\topic}[1]{\par\addvspace{0.5em}\noindent\textbf{#1}\par\nobreak}
\newcommand{\needinfo}[1]{\par\noindent{\color{red!65!black}\textbf{〔待补充〕}#1}\par}
\newcommand{\doi}[1]{\href{https://doi.org/#1}{\nolinkurl{#1}}}
\newcommand{\SE}{\mathrm{SE}(3)}
\newcolumntype{Y}{>{\raggedright\arraybackslash}X}
\newcolumntype{P}[1]{>{\raggedright\arraybackslash}p{#1}}
\renewcommand{\arraystretch}{1.25}
\pagestyle{plain}

\begin{document}

\begin{center}
 {\zihao{3}\bfseries 柔性支撑空间机器人耦合动力学\par
 与学习增强阻抗控制\par}
 \vspace{0.55em}
 {\zihao{-4}航空航天结构力学及控制全国重点实验室开放课题\quad 重点课题}
\end{center}
\vspace{0.35em}

% ==================== 报告正文开始 ====================
\section{立项依据与研究内容}

\subsection{项目的立项依据}

\subsubsection{研究意义}

大型空间天线、空间太阳能电站等装备向大尺度、轻量化发展，使机器人在轨组装成为突破运载包络限制的重要途径。机器人依附于已建桁架搬运和安装模块时，轻质结构同时承担作业支撑与载荷传递功能，其弹性变形直接影响机器人的运动精度。如何在有限结构刚度和不完备动力学模型条件下维持稳定、准确的作业运动，是大型空间结构组装需要解决的基础问题\cite{hu2025,wu2024}。

本项目关注柔性桁架支撑下多自由度机器人的携载运动与接触前停驻。关节驱动产生的反作用载荷激励桁架振动，安装点的位移和转动又沿机械臂传递至作业端，形成运动与结构响应的双向耦合。与此同时，连杆惯性、负载、关节摩擦及执行时延与名义模型存在偏差，使理论控制指令难以在实际机械臂上准确复现。增加反馈刚度虽可减小部分运动误差，却可能提高柔性模态附近的控制作用，放大反作用激励，延长运动结束后的稳定时间。

以作业端为浮动根的倒置动力学建模，可将末端绝对位姿直接引入广义坐标；特殊欧氏群$\SE$上的阻抗控制可统一描述位置与姿态响应；动力学二次规划（QP）则可在驱动约束下实现所需运动。三者为柔性支撑机器人提供了明确的模型控制框架。然而，该框架的控制效果仍受模型失配制约。学习补偿若仅按固定基座条件训练，容易混淆关节内部误差与真实支撑作用；若在优化求解之后直接叠加力矩，又可能改变已规划的阻抗响应和力矩约束。

因此，需要在保持动力学结构和阻抗控制目标的基础上，研究残差强化学习对实际执行误差的补偿机制，明确其与柔性支撑响应、约束优化及不同控制更新频率之间的关系。本项目拟形成适用于柔性支撑机器人的学习增强阻抗控制方法，在相同作业时间和驱动条件下改善末端跟踪与结构振动，为后续接触操作提供稳定的运动控制基础。

\subsubsection{国内外研究现状及发展动态}

\topic{（1）柔性结构与空间机器人的耦合动力学}

空间机器人动力学研究已形成多体建模、动量耦合分析和操作控制等基础方法\cite{papadopoulos2021}。面向大型柔性结构，模态缩减可在保留主要弹性效应的同时降低计算量，Sonneville等对相关方法及其多体应用进行了研究\cite{sonneville2021}。Ni等进一步分析了柔性机械臂动态边界与模态选取的关系，其算例说明，模型精度应结合具体工况和频段判断，适当的定常模态近似也可能满足分析需求\cite{ni2024}。

已有研究为机器人与结构联合建模提供了依据。对实际运控系统，仍需把这些模型与传感、控制坐标和驱动力矩对应起来。尤其在倒置模型中，计算根位于作业端，而物理支撑作用于安装端，两者通过完整惯性耦合相联系。以固定基座模型的参数修正代替支撑动力学，难以保持这一关系；将安装端运动全部视为外部给定扰动，也不足以评价机器人控制对结构振动的反向影响。

\topic{（2）几何阻抗控制与运动振动协调}

阻抗控制通过设计运动误差与作用力之间的动态关系组织机器人响应。Kim等利用$\SE$与其李代数之间的映射，在指数坐标中设计六自由度阻抗，并给出映射微分及其时间导数，为位姿统一控制提供了方法\cite{kim2025}。在柔性支撑机器人研究中，Nenchev等提出反作用零空间控制，利用机械臂运动调节支撑结构振动\cite{nenchev1999}；Patel和Damaren结合柔性坐标估计开展末端与振动控制，验证了柔性状态反馈的作用\cite{patel2024}。

上述方法分别解决了位姿响应设计和柔性耦合控制的重要问题，但两者在受约束实际系统中的结合仍需研究。阻抗给出的参考加速度未必同时满足欠驱动动力学、关节边界和力矩变化率限制。约束激活时，参考跟踪的放松又会改变实际阻抗响应。因此，需要通过动力学优化协调运动要求、结构响应与驱动能力，并分析模型误差下这种协调关系的适用条件。

\topic{（3）模型引导的残差强化学习}

残差强化学习保留已有控制器，只学习其未能充分补偿的控制作用。Johannink等将传统反馈与学习残差结合，在实际装配任务中验证了这一思路\cite{johannink2018}。Li等利用历史数据构造增量系统，研究了高维机器人跟踪的残差学习方法及稳定性问题\cite{li2025}。相关综述进一步指出，学习效果、模型不确定性、物理约束和闭环稳定性需要联合考虑\cite{brunke2022}。

团队近期已开展微分动态规划（DDP）引导的残差强化学习研究，利用名义轨迹、反馈增益及局部价值信息组织学习，通过双端点残差输出连接低频策略与高频力矩控制，并完成Panda机械臂实验。其进一步用于本项目时，面临控制结构的变化：名义控制由给定DDP反馈转为随状态和约束变化的阻抗--QP控制，原有价值函数信息不能直接沿用；柔性支撑又引入训练中未充分覆盖的动态作用。需要重新建立学习状态、奖励和动力学约束之间的联系。

\topic{（4）研究切入点}

现有成果表明，几何阻抗、约束动力学与残差学习具有互补性。本项目拟从\textbf{柔性支撑作用的动力学一致表征，以及残差补偿与受约束阻抗响应的相容性}切入：以倒置建模连接作业端状态和安装端响应，以主导柔性模态描述可观测结构运动，以受约束残差学习减轻实际机械臂执行误差，研究精度改善与振动抑制的共同条件。团队已有模块组装实验和空间多柔体缩比建模成果\cite{shi2024,qi2025}，可为这一研究提供地面验证基础。

\subsubsection{与申请指南的关系}

本项目主要对应指南2.5“闭环系统模型辨识、分析与控制”，面向强耦合、多振动模态及多源不确定性，研究结合深度强化学习的动力学残差表征、闭环分析与运动控制；同时呼应2.2“轻质柔性结构动力学设计与分析”中的在轨组装、刚柔耦合动力学与振动抑制。研究限定于单机器人、既定连接构型和若干代表性支撑状态，以携载运动及接触前停驻构成验证任务，适应两年重点开放课题的实施条件。

\begingroup
\small
\setstretch{1.12}
\renewcommand{\refname}{主要参考文献}
\begin{thebibliography}{99}
\setlength{\itemsep}{0.3em}
\bibitem{hu2025} 胡海岩, 田强, 文浩, 罗凯, 马小飞. 极大空间结构在轨组装的动力学与控制[J]. 力学进展, 2025, 55(1): 1--29. DOI: \doi{10.6052/1000-0992-24-044}.
\bibitem{wu2024} 吴志刚, 蒋建平, 邬树楠, 李庆军, 王兴, 谭述君, 邓子辰. 航天结构空间组装动力学与控制研究进展[J]. 力学进展, 2024, 54(2): 344--390. DOI: \doi{10.6052/1000-0992-23-041}.
\bibitem{papadopoulos2021} Papadopoulos E, Aghili F, Ma O, Lampariello R. Robotic manipulation and capture in space: A survey[J]. Frontiers in Robotics and AI, 2021, 8: 686723. DOI: \doi{10.3389/frobt.2021.686723}.
\bibitem{sonneville2021} Sonneville V, Scapolan M, Shan M, Bauchau O A. Modal reduction procedures for flexible multibody dynamics[J]. Multibody System Dynamics, 2021, 51(4): 377--418. DOI: \doi{10.1007/s11044-020-09770-w}.
\bibitem{ni2024} Ni S, Chen W, Chen T. Dynamic modeling and analysis of multi-flexible-link space manipulators under time-varying dynamic boundary conditions[J]. Advances in Space Research, 2024, 74(10): 5224--5243. DOI: \doi{10.1016/j.asr.2024.07.067}.
\bibitem{kim2025} Kim J, Sung M, Choi Y, Park J, Chung W K. Impedance control design framework using commutative map between $\SE$ and $\mathfrak{se}(3)$[J]. IEEE Transactions on Robotics, 2025, 41: 6193--6212. DOI: \doi{10.1109/TRO.2025.3619066}.
\bibitem{nenchev1999} Nenchev D N, Yoshida K, Vichitkulsawat P, Uchiyama M. Reaction null-space control of flexible structure mounted manipulator systems[J]. IEEE Transactions on Robotics and Automation, 1999, 15(6): 1011--1023. DOI: \doi{10.1109/70.817666}.
\bibitem{patel2024} Patel D, Damaren C J. Tip and vibration control of space robots using estimated flexible coordinates[J]. Frontiers in Space Technologies, 2024, 5: 1447545. DOI: \doi{10.3389/frspt.2024.1447545}.
\bibitem{johannink2018} Johannink T, Bahl S, Nair A, Luo J, Kumar A, Loskyll M, Aparicio Ojea J, Solowjow E, Levine S. Residual reinforcement learning for robot control[EB/OL]. arXiv:1812.03201, 2018. DOI: \doi{10.48550/arXiv.1812.03201}.
\bibitem{li2025} Li C, Liu F, Wang Y, Buss M. Data-informed residual reinforcement learning for high-dimensional robotic tracking control[J]. IEEE/ASME Transactions on Mechatronics, 2025, 30(3): 1681--1691. DOI: \doi{10.1109/TMECH.2024.3412275}.
\bibitem{brunke2022} Brunke L, Greeff M, Hall A W, Yuan Z, Zhou S, Panerati J, Schoellig A P. Safe learning in robotics: From learning-based control to safe reinforcement learning[J]. Annual Review of Control, Robotics, and Autonomous Systems, 2022, 5: 411--444. DOI: \doi{10.1146/annurev-control-042920-020211}.
\bibitem{shi2024} 史玲玲, 肖晓龙, 张晓峰, 范立佳, 单明贺, 田强. 空间机器人在轨组装多模块单元的对接力控制与地面实验[J]. 力学学报, 2024, 56(3): 800--816. DOI: \doi{10.6052/0459-1879-23-458}.
\bibitem{qi2025} 祁伊凡, 单明贺, 田强. 空间多柔体系统缩比等效动力学建模方法研究[J]. 中国科学: 物理学 力学 天文学, 2025, 55(2): 224518. DOI: \doi{10.1360/SSPMA-2024-0302}.
\end{thebibliography}
\endgroup

\subsection{项目的研究内容、研究目标及拟解决的关键科学问题}

\subsubsection{研究内容}

\topic{（1）柔性支撑机器人耦合动力学与受约束阻抗运动控制}

针对安装端弹性运动和末端绝对位姿控制之间的联系，研究以作业端为浮动根的倒置动力学表征，将柔性支撑的主导模态、界面载荷和机器人惯性耦合纳入同一模型。利用位姿、关节和结构响应测量校核模型，分析支撑运动与关节内部失配的作用通道。在此基础上，以$\SE$阻抗生成位姿响应要求，通过动力学QP协调参考运动、结构响应和驱动约束，形成供残差学习依托的名义控制方法。

\topic{（2）多源失配条件下的残差学习补偿与运动振动协调}

针对实际机械臂惯性、摩擦及执行时序偏差，拓展已有DDP引导残差强化学习方法，建立适用于阻抗--QP控制的状态表征、学习目标和补偿接口。利用名义模型、运动预览和历史响应组织有界残差策略，将其输出纳入总力矩约束，研究不同更新频率、约束激活及柔性模态对学习补偿效果的影响。以误差反馈和振动响应共同约束补偿作用，分析闭环稳定条件，通过柔性桁架地面实验验证方法有效性。

两项内容依次解决名义动力学与控制的建立、实际模型误差的补偿。实验围绕上述方法展开，形成由模型分析、算法设计到真实弹性支撑验证的完整研究过程。

\subsubsection{研究目标与考核指标}

总体目标是揭示柔性支撑作用、机械臂模型失配与末端阻抗响应之间的关系，建立倒置动力学约束下的残差学习控制方法，改善携载运动精度和停驻振动，并明确其在代表性支撑及载荷条件下的适用范围。

拟在现有七自由度Franka平台上构建3种可重复支撑状态和2种载荷条件，采用表\ref{tab:targets}评价所提方法。各项数值为项目拟达到的目标。

\begin{table}[htbp]
\centering
\caption{拟定技术指标与评价条件}
\label{tab:targets}
\small
\begin{tabularx}{\linewidth}{P{2.3cm}YP{4.1cm}}
\toprule
评价内容 & 拟达到的指标 & 评价条件 \\
\midrule
运动精度 & 末端位置跟踪均方根误差较名义阻抗--QP控制降低不小于20\%，姿态均方根误差不增加。 & 在相同任务时间和驱动约束下，覆盖载荷变化及主要惯性参数$\pm10\%$建模偏差。 \\
结构振动 & 运动结束后固定观察窗内，支撑测点动态位移均方根值较上述基线降低不小于25\%。 & 与运动精度在同组实验中评价，采用真实柔性支撑在停驻阶段的衰减响应。 \\
多频率执行 & 指令力矩变化率均方根值较零阶保持残差接口降低不小于15\%，位置误差不增加。 & 比较相同训练预算、控制频率及物理约束下的两种残差接口。 \\
多工况验证 & 完成6组支撑与载荷组合，每组覆盖2类轨迹，每类不少于5次配对重复试验。 & 包含未参与策略训练与调参的轨迹；逐工况报告误差、振动及约束激活情况。 \\
\bottomrule
\end{tabularx}
\end{table}

名义基线与所提方法采用同一倒置模型、状态估计、阻抗参数和QP权重，差别为是否引入学习残差；另设置模型摩擦补偿及现有监督学习残差作为对照。训练与评价工况预先划分，基线在相同训练工况充分整定。指标按各预定工况重复实验的均值分别评价，同时给出离散程度，保留未达目标和触发约束的结果。

\subsubsection{拟解决的关键科学问题}

\topic{（1）柔性支撑与机械臂内部失配如何共同改变末端阻抗响应}

支撑运动和关节内部误差均可表现为末端偏差，但前者是安装界面的真实动力学作用，后者包含驱动能够补偿的执行误差，两者的作用空间不同。拟研究倒置坐标下界面载荷、柔性模态和关节残差的传递关系，确定有限测量条件下可分辨的误差成分及驱动可补偿范围，避免将真实结构响应吸收为应被抵消的内部残差。

\topic{（2）受约束、多更新频率的残差补偿如何兼顾精度与振动耗散}

残差策略改变动力学QP的力矩分配，约束激活又改变策略作用到实际运动的映射；低频策略与高频控制之间的时序差异可能向柔性模态注入额外激励。拟研究补偿幅值、变化率、反馈带宽与闭环能量变化之间的关系，建立学习残差与阻抗响应相容的条件，明确提高跟踪精度同时抑制结构振动所需的控制边界。

\subsection{拟采取的研究方案及可行性分析}

\subsubsection{总体研究思路与技术路线}

以现有倒置建模和阻抗--QP控制为起点，补充柔性安装界面的动态载荷描述；继承残差强化学习中的模型引导、时序记忆和双端点输出，重新组织与QP相容的学习目标；通过整体模型分析、分层对照和地面实验逐步验证。技术路线见图\ref{fig:route}。

\begin{figure}[htbp]
\centering
\begin{tikzpicture}[
 box/.style={draw=black!70,rounded corners=2pt,align=center,inner sep=7pt,font=\small,text width=5.6cm,minimum height=1.1cm},
 wide/.style={box,text width=11.6cm},
 arr/.style={-{Latex[length=2mm]},thick,draw=black!75}]
 \node[wide] (task) at (0,0) {\textbf{柔性桁架支撑下的机器人携载运动与停驻}\\末端绝对位姿\quad 安装界面响应\quad 模型与执行偏差};
 \node[box] (model) at (-3.2,-1.8) {\textbf{倒置动力学与主导柔性模态}\\坐标及功率一致\quad 动态支撑载荷};
 \node[box] (learn) at (3.2,-1.8) {\textbf{模型引导的残差强化学习}\\时序特征\quad 多源失配训练};
 \node[box] (imp) at (-3.2,-3.6) {\textbf{$\SE$阻抗运动要求}\\位姿统一表征\quad 参考加速度};
 \node[box] (res) at (3.2,-3.6) {\textbf{有界双端点残差}\\低频策略输出\quad 高频平滑执行};
 \node[wide] (qp) at (0,-5.4) {\textbf{统一动力学QP}\quad 残差计入总力矩约束\\运动跟踪\quad 柔性响应调节\quad 力矩及变化率限制};
 \node[wide] (test) at (0,-7.2) {\textbf{闭环分析与地面验证}\\精度与振动共同评价\quad 多工况对照\quad 适用范围};
 \draw[arr] (task.south) -- ++(0,-0.25) -| (model.north);
 \draw[arr] (task.south) -- ++(0,-0.25) -| (learn.north);
 \draw[arr] (model) -- (imp);
 \draw[arr] (learn) -- (res);
 \draw[arr] (imp.south) -- ++(0,-0.25) -| ($(qp.north)+(-2.3,0)$);
 \draw[arr] (res.south) -- ++(0,-0.25) -| ($(qp.north)+(2.3,0)$);
 \draw[arr] (qp) -- (test);
\end{tikzpicture}
\caption{项目技术路线}
\label{fig:route}
\end{figure}
\FloatBarrier

\subsubsection{倒置动力学与柔性支撑作用的统一表征}

取作业端位姿$T_E\in\SE$与反向排列的关节坐标$\bm q_j$描述机器人构型，广义速度为$\bm\nu=[\bm V_E^{\mathsf T},\dot{\bm q}_j^{\mathsf T}]^{\mathsf T}$。倒置建模改变动力学树的根及递推方向，物理安装端仍与桁架相连。根据虚功一致性推导速度、惯性矩阵和广义力的变换，校核正向模型与倒置模型在相同物理运动下的能量和驱动力矩对应关系。

在小弹性变形的代表性支撑状态下，取少量主导模态$\bm\eta$描述结构响应。名义耦合关系写为
\begin{align}
 \widehat{\bm M}_r\dot{\bm\nu}+\widehat{\bm h}_r
 &=\bm S^{\mathsf T}(\bm z-\bm d)
   +\bm J_b^{\mathsf T}\bm w_b, \label{eq:robot}\\
 \bm M_\eta\ddot{\bm\eta}+\bm C_\eta\dot{\bm\eta}+\bm K_\eta\bm\eta
 &=-\bm B_\eta^{\mathsf T}\bm w_b+\bm f_\eta. \label{eq:modal}
\end{align}
其中，$\bm w_b$为支撑对机器人的界面力和力矩，$\bm J_b$和$\bm B_\eta$将两侧速度映射到同一界面坐标；联立安装界面运动相容条件，保持作用与反作用的功率一致。$\bm z$为包含补偿的等效总关节指令，$\bm d$为命令侧残差补偿，$\bm S$为驱动选择矩阵。实际系统与该名义关系的偏差单独保留，用于误差分析。

融合动捕、IMU和编码器获得作业端及安装端运动，通过结构测点校核主导模态与末端误差的关系。闭环数据用于确定敏感支撑和载荷参数，以独立轨迹校核预测。关节残差描述驱动通道内的等效执行偏差，六维支撑作用及驱动未覆盖的误差由结构模型、状态估计和鲁棒反馈处理。

\subsubsection{面向结构响应的阻抗--QP控制}

以相对位姿$T_E^{-1}T_d$的对数坐标$\bm e$统一表示位姿误差，结合指数映射微分，将设计的虚拟惯量、阻尼和刚度转化为作业端参考加速度$\bm a_E^{\rm imp}$。平移和转动采用相应物理尺度归一化；阻抗参数结合支撑主导频率选取，研究运动频段与振动频段的反馈作用分配。

在当前构型处，将机器人动力学、模态方程及界面加速度相容条件联立，构建以广义加速度$\bm a=\dot{\bm\nu}$、模态加速度、界面载荷和总指令$\bm z$为变量的QP。目标函数拟采用
\begin{equation}
 \frac12\|\bm a_E-\bm a_E^{\rm imp}\|_{\bm W_E}^2
 +\frac12\|\bm a_j-\bm a_j^{\rm ref}\|_{\bm W_j}^2
 +\frac12\|\dot{\bm\eta}+\Delta t\ddot{\bm\eta}\|_{\bm W_\eta}^2
 +\frac12\|\bm z\|_{\bm R}^2 ,
 \label{eq:qp}
\end{equation}
其中$\|\bm x\|_{\bm W}^2=\bm x^{\mathsf T}\bm W\bm x$，$\bm a_j^{\rm ref}$提供关节运动正则。动力学等式和关节、力矩及变化率边界作为约束，末端加速度采用软跟踪，利用结构响应项调节有限驱动能力下的运动与振动关系。

残差$\bm d$在求解前确定，通过式\eqref{eq:robot}进入动力学。总力矩及其变化约束直接作用于$\bm z$，并与实际执行接口的重力补偿约定和上一拍指令回显对应。由此，学习补偿占用的驱动余量在求解时即可被计入。研究约束激活、界面模型误差和模态截断对跟踪松弛量的影响，确定权重与运动带宽的适用范围。

\subsubsection{与阻抗--QP相容的残差强化学习}

\topic{（1）建立包含结构响应和控制意图的学习状态}

在已有状态、历史和参考预览编码的基础上，加入安装端运动、可观测柔性响应及前一成功QP的约束状态。期望速度和加速度来自任务轨迹，模型力矩和惯性耦合量提供物理特征；历史窗口描述摩擦、执行滞后及观测时序的影响。策略使用当前可获得的测量和既往求解结果，避免由本拍未知QP解构造本拍动作。现有监督学习GRU可提供对照和特征设计经验，残差强化学习策略按闭环任务重新训练。

\topic{（2）拓展模型引导的训练目标}

原DDP方法中的局部价值信息依赖既定优化轨迹。本项目拟利用阻抗误差储能及局部动力学响应灵敏度构造学习权重，将位置姿态误差、主导模态能量、残差幅值与变化率共同纳入奖励；在固定约束激活模式下分析局部灵敏度，并检验切换时的变化。奖励设置以相同运动时间内的精度和振动共同改善为目标。

训练采用近端策略优化，结合随机化模型的逆动力学构造辅助残差目标，并保留相邻策略区间的动作一致性约束。辅助目标只用于训练中具有可信动力学信息的起始状态，未来状态不作为在线可用量。随机化覆盖连杆惯性、负载、摩擦、支撑刚度与阻尼以及观测和执行时序；参数范围由标定与预实验确定。训练以耦合仿真为主，利用实测数据校核模型及分布差异，真机执行采用冻结策略在线推理，后续更新在试验批次之间完成。

\topic{（3）实现低频学习与高频约束控制的衔接}

策略按低频时刻$t_j$输出区间两端残差，经滤波和幅值处理后，在高频控制周期内插值得到
\begin{equation}
 \bm d(t)=(1-\beta)\bar{\bm d}^{\,0}_j+\beta\bar{\bm d}^{\,1}_j,
 \qquad \beta=\frac{t-t_j}{T_\pi},\quad t_j\le t<t_j+T_\pi.
 \label{eq:residual}
\end{equation}
以现有100~Hz策略和1~kHz力矩控制为初始配置，逐周期把插值结果送入QP，并限制实际总指令的变化。相邻区间衔接、状态过期和推理缺帧均在训练及验证中覆盖，记录策略延迟、QP计算时间和指令变化。处理新旧残差、摩擦模型与动量观测补偿之间的关系，先统一各自估计的物理量和力矩来源，再确定组合方式。

\topic{（4）分析补偿与整体闭环响应的相容条件}

构造包含位姿误差与柔性模态能量的候选Lyapunov函数，分析残差误差、状态估计误差、加速度跟踪松弛和采样时序对其变化率的影响。在名义控制局部稳定、QP可行及误差有界的条件下，推导闭环误差最终有界的充分条件，给出补偿增益、变化率和容许时延的限制。对难以满足能量变化条件的补偿方向降低作用强度，并检验收缩补偿后QP的可行性；分析同时覆盖约束切换与未保留模态，不将指令限幅本身作为稳定性结论。

\subsubsection{地面实验方案与评价方法}

采用已有Franka机械臂、力矩控制接口、动捕和IMU，增配可调柔性桁架支撑及必要测点。验证分为固定基座算法复核、可重复安装端运动测试、真实柔性支撑闭环实验三个层次。前两层校核坐标、时序及策略迁移，第三层使机器人驱动力真实作用于支撑结构，用于评价运动与振动的双向耦合。

设置平滑携载转运和连续轨迹运动两类任务，均包含终点停驻。3种支撑状态与2种载荷构成6组工况，试验中保持连接关系固定。各方法使用相同参考、执行时间、传感配置和硬约束，成对重复并交替安排顺序。依次比较名义阻抗--QP、模型摩擦补偿、监督学习残差和本项目方法；再以去除结构响应特征、振动奖励或双端点接口的消融，辨别性能改善来源。

末端精度用位置误差和旋转测地误差的均方根值评价：
\begin{equation}
 E_p=\sqrt{\frac1N\sum_{k=1}^N\|\bm p_k-\bm p_{d,k}\|^2},\qquad
 E_R=\sqrt{\frac1N\sum_{k=1}^N
 \|\operatorname{Log}(\bm R_{d,k}^{\mathsf T}\bm R_k)^\vee\|^2}.
 \label{eq:metric}
\end{equation}
停驻振动以相对静态平衡位置的测点动态位移均方根值评价，观察窗取该支撑状态一阶固有周期的5倍，各方法共用同一窗口。同步报告力矩、变化率、QP松弛和约束激活比例。外部原始测量按时间戳去重并与参考对齐，区分控制估计误差与测量评价误差；信号接近噪声底时报告绝对值。地面重力、附加支承阻尼和模态相似性单独标定，明确验证结论的适用范围。

\subsubsection{可行性分析}

本项目已有可直接衔接的算法和实验基础。现有Franka控制程序具备倒置动力学、$\SE$参考加速度、动力学QP以及命令侧残差接口；DDP引导残差学习已完成固定基座真机验证，提供了时序编码、训练和多更新频率部署经验。团队的柔性多体建模及地面平台可支持支撑结构建模和标定。新增研究集中于动态支撑、学习目标转换与整体闭环分析，能够沿现有系统逐步实施。

对支撑与内部残差混淆问题，采用独立结构测点和对照激励检验；对迁移不足，依据未参与训练的数据调整仿真分布，并在相同条件下与监督学习残差比较；对实时计算压力，保留少量主导模态，将网络推理与高频QP分频执行；对振动放大，结合模态测试收缩补偿带宽与增益。每阶段均以可观测响应和对照试验判定改进效果。

\subsection{本项目的特色与创新之处}

\topic{（1）面向柔性安装界面建立动力学一致的阻抗控制方法}

在已有倒置建模基础上，将作业端位姿、安装界面动态载荷及主导柔性模态关联起来，以同一动力学约束协调阻抗运动要求和结构响应。拟阐明支撑作用与内部失配的不同传递途径，形成适用于柔性支撑机器人的控制表征和约束分配方法。

\topic{（2）形成兼顾阻抗响应与结构振动的模型引导残差学习方法}

将已有DDP引导残差学习拓展至随状态和约束变化的阻抗--QP系统，利用物理响应和任务误差组织学习，将有界残差纳入统一力矩预算，并结合双端点执行研究不同更新频率下的振动激励。拟建立学习补偿与名义阻抗的相容条件，使模型误差补偿同时服务于末端精度和结构响应改善。

\subsection{年度研究计划及预期研究成果}

\subsubsection{年度研究计划}

\topic{2027年1月至6月}
确定支撑与载荷工况，完成柔性试件设计和模态标定；建立倒置模型、安装界面及主导模态的耦合关系，复核现有控制和残差接口。阶段成果为耦合模型、标定数据及基线实验方案。

\topic{2027年7月至12月}
完成受约束阻抗控制及初步闭环分析，建立包含结构响应的残差训练环境；研究状态、奖励及辅助监督设计，完成固定基座与安装端运动条件下的对照。形成论文初稿1篇、年度进展报告，并与重点实验室开展专题交流。

\topic{2028年1月至6月}
完成残差策略、双端点接口与动力学QP融合，完善约束切换和柔性模态影响分析；开展真实柔性支撑预实验、迁移校核与消融比较，形成方法论文和发明专利申请材料。

\topic{2028年7月至12月}
完成6组工况的配对重复实验和独立评价，明确算法适用范围；整理程序、数据及测试报告，完成论文发表或录用、专利申请和结题材料，与重点实验室开展成果交流。

\subsubsection{预期研究成果与学术交流}

形成柔性支撑机器人耦合动力学与阻抗--QP程序1套，模型引导残差学习控制算法1套，配套多工况实验数据、评测脚本及使用说明；完成地面原理验证，提交技术总结报告和实验测试报告各1份。围绕新增研究发表或录用学术论文2篇，申请发明专利1项，培养参与研究的研究生2名。

每年至少与重点实验室研究人员开展1次实质性学术交流，结合模型、实验和阶段成果组织研讨与学术报告；拟参加1次相关国内学术会议，并利用线上形式开展国际学术讨论。项目成果署名、资助标注和知识产权安排按重点实验室开放课题管理要求执行。

\section{研究基础与工作条件}

\subsection{工作基础}

申请人所在团队长期从事空间在轨操作动力学与控制研究，在柔性多体建模、Franka机械臂运控、残差学习和模块对接方面形成了连续积累。与本项目直接相关的基础如下。

\topic{（1）柔性多体动力学与地面模拟基础}

申请人参与了柔性多体系统模态缩减研究，相关成果发表于\textit{Multibody System Dynamics}\cite{sonneville2021}；团队进一步开展空间多柔体系统缩比等效建模，研究相似关系及参数灵敏度对缩比模型设计的影响\cite{qi2025}。现有约$4\,\mathrm m\times10\,\mathrm m$的在轨操作微重力模拟地面平台，以及已开展的多模块对接力控制实验\cite{shi2024}，为本项目的模型校核、支撑试验和数据评价提供了基础。

\topic{（2）倒置动力学、几何阻抗与QP控制实现}

团队已围绕七自由度Franka FR3开展运动控制系统建设，形成以作业端侧link7为浮动根、沿关节链反向组织连杆的动力学实现。模型包含六维根运动和七维关节运动，实际安装端link0的支撑作用通过广义力进入动力学。现有程序结合动捕、IMU和关节测量形成状态输入，以$\SE$阻抗生成六维参考加速度，通过动力学QP求取关节加速度与力矩，并设置关节、力矩及其变化边界。

系统已具备位姿保持、圆轨迹和固定基座轨迹实验入口，以及记录、回放和接口校核流程。支撑基线目前采用静态载荷近似，GRU支路主要使用固定基座数据。动态支撑描述与柔性工况的闭环验证将由本项目继续开展。

\topic{（3）DDP引导残差强化学习与真机验证}

团队近期完成了DDP引导的力矩残差强化学习研究，相关稿件已提交IEEE Robotics and Automation Letters。方法保留名义前馈与反馈，用DDP轨迹和局部价值信息引导学习，通过当前状态、历史及预览编码生成有界残差，结合逆动力学辅助监督和动作一致性训练，以双端点输出连接100~Hz策略与1~kHz力矩控制。

稿件已报告Panda在训练轨迹、未参与训练轨迹及未建模负载条件下的实验。表\ref{tab:preliminary}摘录未参与训练轨迹集合的结果。末端指标沿用稿件的平均误差，关节指标为均方根误差，与本项目的末端均方根考核分别评价。

\begin{table}[htbp]
\centering
\caption{已有Panda实验结果：未参与训练的轨迹集合}
\label{tab:preliminary}
\small
\begin{tabularx}{\linewidth}{Yccc}
\toprule
方法 & 末端平均位置误差/mm & 平均姿态误差/(${}^\circ$) & 关节位置RMS/(${}^\circ$) \\
\midrule
名义DDP & 2.565 & 1.348 & 1.216 \\
摩擦补偿DDP & 1.962 & 1.184 & 0.968 \\
DDP引导残差学习 & 0.338 & 0.203 & 0.110 \\
\bottomrule
\end{tabularx}
\vspace{0.35em}
\begin{minipage}{\linewidth}
\footnotesize 来源：团队2026年投稿稿件，表V，Held-out集合。该结果对应固定安装的Panda及DDP控制；柔性支撑FR3上的阻抗--QP融合尚需本项目验证。
\end{minipage}
\end{table}

稿件的未建模负载及零阶保持接口对比，为策略迁移和多更新频率执行提供了经验。柔性支撑条件下的阻抗--QP策略仍需重新训练与验证。

\topic{（4）动力学残差辨识与补偿接口基础}

团队已开展基于标定动力学、动量滤波特征和GRU的监督学习残差研究，形成关节及电机状态、规划参考和模型量的特征构造与推理实现。现有QP接口将命令侧残差计入总力矩变量，并区分该残差与物理外部广义力，为本项目建立动力学一致的补偿接口提供了直接经验。

上述GRU采用监督学习，与强化学习策略分别训练。固定基座特征向动态支撑条件的迁移，以及命令侧残差与测量侧动量观测的统一，已开展方案分析，需进一步通过闭环实验验证。

\topic{（5）接触几何与柔顺对接的相关研究积累}

团队近期研究稿围绕复杂接口对接，提出物理方向辅助接触场配准和配准锚定捕获：利用接触位姿及物理方向信息构造接触场，通过在线接触位姿估计目标，再以配准结果为锚点进行有界参考修正。该研究建立了参考更新子系统的离散储能关系，已开展仿真及真机接触记录配准回放，新方法的真机捕获实验仍待完成。

这一工作说明，后续接触配准和参考修正需要下层控制准确执行给定的位姿及柔顺响应。本项目据此明确末端参考接口和响应评价需求，重点解决柔性支撑条件下的运动控制基础问题。既有接触研究作为应用与方法积累，为后续向完整组装对接任务发展提供衔接。

\subsection{工作条件}

\topic{（1）已具备的条件}

依托单位具有空间结构动力学和机器人控制研究团队，已具备七自由度Franka机械臂、开放力矩控制程序、动捕与IMU输入及同步记录基础。固定基座Panda残差学习实验提供了仿真训练到真实力矩控制的经验，FR3倒置建模与阻抗--QP程序提供了具体集成基础。现有地面模拟平台及模块对接实验条件可支持控制验证、场地布置和专项试验组织。

\topic{（2）尚需完善的条件及解决途径}

拟根据现有机器人承载能力增配可调柔性桁架支撑、标准载荷和结构测点，完成模态、时间同步及空间外参标定；共享现有计算与测量设备，补充必要的同步采集和控制接口。重点完善真实弹性反馈、安装界面响应和独立末端测量，支撑试件、接口改装及专项测试纳入项目预算。

拟与重点实验室固定研究人员围绕结构模态测试、耦合动力学分析和学习增强控制开展合作，结合双方实际条件安排模型讨论、联合验证及成果交流。具体设备共享及试验分工依据实际合作安排落实。

\needinfo{南航重点实验室合作人员姓名、职称及分工；现有动捕、IMU和结构测量设备的型号、精度及可共享条件。}

\subsection{申请人简介}

\topic{（1）学历与研究工作简历}

单明贺，北京理工大学教授、博士生导师，主要从事空间在轨操作动力学与控制、柔性多体动力学和空间机器人技术研究。2007年9月至2011年7月在哈尔滨工业大学学习，获飞行器制造工程学士学位；2011年9月至2013年7月在哈尔滨工业大学学习，获航空宇航制造工程硕士学位；2013年11月至2018年6月在荷兰代尔夫特理工大学学习，获航天工程博士学位。2018年11月至2020年11月在美国马里兰大学帕克分校从事博士后研究，2021年1月起在北京理工大学工作。

申请人主持国家自然科学基金等科研任务，在空间在轨服务及空间结构动力学相关方向形成了持续研究积累。其研究经历与本项目的柔性多体动力学、机器人运动控制及地面验证紧密相关。

\topic{（2）近期与本项目有关的主要论著}

\begin{enumerate}[label=\arabic*.]
 \begin{samepage}
 \item 祁伊凡, 单明贺, 田强. 空间多柔体系统缩比等效动力学建模方法研究[J]. 中国科学: 物理学 力学 天文学, 2025, 55(2): 224518. DOI: \doi{10.1360/SSPMA-2024-0302}.
 \end{samepage}
 \begin{samepage}
 \item 史玲玲, 肖晓龙, 张晓峰, 范立佳, 单明贺, 田强. 空间机器人在轨组装多模块单元的对接力控制与地面实验[J]. 力学学报, 2024, 56(3): 800--816. DOI: \doi{10.6052/0459-1879-23-458}.
 \end{samepage}
 \begin{samepage}
 \item Sonneville V, Scapolan M, Shan M, Bauchau O A. Modal reduction procedures for flexible multibody dynamics[J]. Multibody System Dynamics, 2021, 51(4): 377--418. DOI: \doi{10.1007/s11044-020-09770-w}.
 \end{samepage}
\end{enumerate}

% 两篇新增文件为匿名投稿稿/研究稿，作者名单未提供，已在“工作基础”说明。
% 不将匿名稿件冒列为申请人已发表论著，正式论著清单按实际发表状态更新。
\topic{（3）学术奖励}

单明贺，国际空间研究委员会青年科学家最佳论文奖（COSPAR Outstanding Paper Award for Young Scientists），2022年。

\needinfo{获奖论文及证书载明信息；中国振动工程学会科学技术二等奖的项目名称、授奖年份、全部受奖人员及本人排序。}

\subsection{项目组主要成员简介}

项目拟按耦合动力学与名义控制、残差学习与分析、实验验证配置人员。申请人负责总体方案、关键方法和项目组织，依托单位研究骨干及研究生承担模型实现、策略训练和实验工作，重点实验室合作人员参与结构动力学分析、试验方案论证及学术交流。

\needinfo{按实际参研名单逐人列明简介、具体任务、作者及出版信息完整的近期论著，以及受奖人员和授奖信息完整的学术奖励。}

\subsection{近五年承担科研项目情况}

既有空间结构动力学、Franka运动控制、残差学习及接触对接研究，为本项目提供理论、算法与实验基础。本项目新增任务集中于柔性安装界面的动力学描述，以及残差强化学习与受约束阻抗控制的融合，交付对象为相应模型、算法和柔性支撑专项实验结果。

\needinfo{近五年各科研项目的正式名称、编号、经费来源、起止年月和承担角色，并依据实际任务书逐项说明与本项目的联系、任务边界及成果增量。}

\clearpage
\section{经费预算}

% 重点课题20--30万元，本稿按30万元编制；数量和单价为需求测算，非供应商报价。
申请经费30.00万元，围绕柔性桁架支撑、测量与控制接口、专项测试及成果产出安排。依托现有机器人、计算设备和地面实验平台开展研究，优先共享可用仪器设备。经费预算如下。

\begingroup
\small
\begin{longtable}{P{3.15cm}>{\centering\arraybackslash}p{1.35cm}P{9.6cm}}
\toprule
预算科目 & 金额（万元） & 计算根据及理由 \\
\midrule
\endfirsthead
\toprule
预算科目 & 金额（万元） & 计算根据及理由 \\
\midrule
\endhead
\midrule
\multicolumn{3}{r}{续下页}\\
\endfoot
\bottomrule
\endlastfoot
\textbf{合计} & \textbf{30.00} & 服务于两项研究内容及其地面实验验证。 \\
设备费 & 4.00 & 同步采集和实时控制接口改装所需设备及配件1套，按4.00万元测算，用于关节、结构响应与末端测量的时序同步。 \\
材料费 & 6.00 & 3组柔性支撑状态对应的桁架及可换连接组件原材料，每组1.20万元，共3.60万元；夹具及标准载荷材料1.40万元；传感安装、线缆等耗材1.00万元。 \\
测试化验加工费 & 9.00 & 支撑及夹具加工装配3批，每批1.50万元，共4.50万元；专项模态与测量标定3批，每批0.80万元，共2.40万元；外协动态测量与联合测试3批，每批0.70万元，共2.10万元。共享仪器免费使用部分不重复计费。 \\
差旅费 & 3.00 & 合作试验、专题交流及参加国内学术会议等约10人次，按每人次0.30万元测算。 \\
会议费 & 0.50 & 项目专题研讨及实验方案评议2次，按每次0.25万元测算。 \\
国际合作与交流费 & 0.00 & 以线上学术交流为主，不安排出国经费。 \\
出版/文献/信息传播/知识产权事务费 & 3.50 & 论文出版相关费用2篇，每篇按1.00万元测算，共2.00万元；发明专利申请及相关事务1.00万元；文献和信息资料0.50万元。 \\
专家咨询费 & 1.00 & 动力学建模及实验方案咨询约5人次，按每人次0.20万元测算；占总经费3.33\%。 \\
劳务费 & 3.00 & 参与建模、策略训练和实验工作的研究生2人，每人累计15个月，每人每月0.10万元；占总经费10.00\%。 \\
\end{longtable}
\endgroup

预算按研究需求测算，设备与测试项目结合现有条件核定，相关支出按实际发生及适用标准执行。专家咨询费和劳务费分别满足申请书模板规定的不超过总经费5\%和10\%的比例要求。

% ==================== 报告正文结束 ====================
\end{document}
